\begin{textsample}{POS Dim 3 – 2021\_c5 – Score 2.00 – 2021\_c5/enade\_2021\_art\_vis\_426539035c.txt}  \label{ex:f3_pos_088}
CAPÍTULO 5 DISTRIBUIÇÃO DOS CONCEITOS 5.1 PANORAMA NACIONAL DA DISTRIBUIÇÃO DOS CONCEITOS O Conceito Enade é um indicador de qualidade que avalia os cursos por intermédio dos desempenhos dos estudantes a partir dos resultados do Enade.
Os valores possíveis vão de 1 ( pior situação ) a 5 ( melhor situação ).
Os cursos com apenas um ou sem nenhum concluinte participante não obtêm o Conceito Enade, ficando Sem Conceito ( SC ).
Enfatiza -se, mais uma vez, que a diferença entre os cursos tabulados neste capítulo e no Capítulo 2 corresponde aos cursos que não foram avaliados, ou seja, àqueles, em princípio, sem alunos concluintes que tivessem sido inscritos para a prova.
Na Tabela 5.1, são apresentadas a quantidade e a distribuição de cursos de Artes Visuais ( Licenciatura ) participantes do Enade/2021, por faixa de conceito e Grande Região.
Observando -se os dados da Tabela 5.1, nota -se que, dos 85 cursos participantes, 29 ( 34,1 \% ) classificaram -se com conceito 4, o valor modal.
Esse foi também o conceito modal em três das cinco regiões : Nordeste ( 46,7 \% ), Sul ( 43,5 \% ) e Centro-Oeste ( 37,5 \% ).
O conceito 3 foi o segundo mais frequente em nível nacional ( 28,2 \%, correspondendo a 24 cursos ), sendo o conceito modal da região Sudeste ( 30,3 \% ), enquanto o conceito 2 foi o terceiro ( 23,5 \%, correspondendo a 20 cursos ).
Houve, ainda, oito cursos ( 9,4 \% ) que receberam conceito 5, e três cursos ( 3,5 \% ) que receberam conceito 1.
Dos 85 cursos de Artes Visuais ( Licenciatura ), um ( 1,2 \% ) ficou Sem Conceito ( SC ).
A região Norte participou com seis cursos ou 7,1 \% do total nacional.
Desses, quatro cursos receberam os conceitos 1 e 2, o que equivale a 33,3 \% do total regional para cada conceito, tendo sido esses os conceitos modais na região.
Os outros dois cursos receberam, cada um, os conceitos 3 e 4, o que equivale a 16,7 \% do total regional.
Nenhum curso ficou Sem Conceito ( SC ) ou recebeu os conceitos 1 e 5, nessa região.
A região Nordeste participou com 15 cursos ou 17,6 \% do total nacional.
Nessa região, sete cursos ( 46,7 \% em termos regionais ) obtiveram conceito 4, o conceito modal para a região, como já comentado.
Os conceitos 2 e 3 foram atribuídos a oito cursos, sendo quatro ( 26,7 \% ), cada.
Ainda nessa região, nenhum curso ficou Sem Conceito ( SC ) ou recebeu os conceitos 1 e 5.
Dos 33 cursos participantes da região Sudeste ( 38,8 \% do total nacional ), o conceito 3 foi o modal, como já anteriormente assinalado, tendo sido obtido por dez cursos ( 30,3 \% ).
O conceito 1 foi atribuído a um curso ( 3,0 \% ), e o conceito 2, a sete cursos ( 21,2 \% ).
Receberam os conceitos 4 e 5, respectivamente, oito cursos ( 24,2 \% ) e sete cursos ( 21,2 \% ).
Nessa região, nenhum curso ficou Sem Conceito ( SC ).
Os 23 cursos da região Sul corresponderam a 27,1 \% do total nacional.
Como já indicado, houve predominância do conceito 4, atribuído a 43,5 \% dos cursos dessa região ( conceito modal ), ou seja, atribuído a dez dos 23 cursos participantes.
Sete cursos obtiveram o conceito 3 ( 30,4 \% ).
O conceito 2 foi atribuído a quatro cursos ( 17,4 \% ), e o conceito 5, a apenas um curso ( 4,3 \% ).
Nenhum curso recebeu conceito 1 e, apenas, um curso ficou Sem Conceito ( SC ).
Já os oito cursos participantes na região Centro-Oeste corresponderam a 9,4 \% do total nacional.
Os conceitos 2 e 4 foram atribuídos a três cursos, cada ( 37,5 \% ), correspondendo a uma distribuição bimodal.
Dois cursos ( 25,0 \% ) obtiveram o conceito 3.
Nessa região, nenhum curso ficou Sem Conceito ( SC ) ou recebeu os conceitos 1 ou 5.
No Gráfico 5.1, é apresentada a distribuição cumulativa do Conceito Enade, segundo a Grande Região.
Em um gráfico de distribuição cumulativa, em cada valor de abscissa ( eixo x ), o valor apresentado no eixo das ordenadas ( eixo y ) é o valor acumulado da variável ( Conceito Enade ).
Observando -se a linha vertical pontilhada, que \textbf{ajuda} a visualizar o percentual de conceitos acumulados até o valor 3, ou seja, conceitos 1, 2 e 3, pode -se notar que, por exemplo, a região Norte ( linha azul ) apresenta quase 85 \% dos cursos com conceito 3 ou abaixo desse valor.
A região Sul ( linha rosa ), por outro lado, apresenta 50 \% dos cursos na mesma situação : com conceito 3 ou abaixo.
Destaca -se que, em algumas regiões, o conceito 5 pode não ter sido alcançado por nenhum curso.
É o que ocorre, por exemplo, nas regiões Norte ( linha azul ), Nordeste ( linha verde ) e Centro-Oeste ( linha \textbf{laranja} ), em que todos os cursos têm conceito menor ou igual a 4, e a poligonal alcança 100 \% no conceito 4.
Em linhas gerais, em um gráfico de distribuição cumulativa, poligonais mais à esquerda representam Grandes Regiões com pior distribuição, e poligonais mais à direita, Grandes Regiões com melhores conceitos.
A região Norte ( linha azul ) e a Centro-Oeste ( linha \textbf{laranja} ) apresentam os cursos com conceitos mais baixos, ao passo que as regiões Sudeste ( linha vermelha ) e Sul ( linha rosa ) apresentam a distribuição com valores maiores.
Cursos em IES situadas na região Nordeste ( linha verde ) apresentam uma situação intermediária. 5.2 CONCEITOS POR CATEGORIA ADMINISTRATIVA, POR MODALIDADE DE ENSINO E GRANDE REGIÃO Na Tabela 5.2, é apresentada a distribuição dos cursos de Artes Visuais ( Licenciatura ) participantes do Enade/2021, por Categoria Administrativa e Modalidade de Ensino, de acordo com os conceitos por eles alcançados, segundo as Grandes Regiões brasileiras.
Dos 85 cursos participantes, 56 ( 65,9 \% ) eram ministrados em Instituições Públicas, e 29 ( 34,1 \% ), em Instituições Privadas.
Quanto à Modalidade de Ensino, 20 cursos eram a Distância, e 65, Presenciais.
De acordo com as informações da Tabela 5.2, em termos nacionais, entre Instituições Públicas e Privadas, observa -se que, dos oito cursos avaliados com conceito 5, sete eram oferecidos em IES Públicas, e apenas um curso, em IES Privadas.
Dos 56 cursos participantes de IES Públicas, o conceito 4 foi o valor modal, atribuído a 25 cursos ( 44,6 \% ).
Dos demais cursos de Instituições Públicas participantes, 12 cursos ( 21,4 \% ) receberam conceito 3, nove cursos ( 16,1 \% ) receberam conceito 2, sete cursos ( 12,5 \% ), conceito 5, e dois cursos ( 3,6 \% ), conceito 1.
Nessa categoria, apenas um curso ficou Sem Conceito ( SC ).
Na rede privada, o conceito modal foi 3, atribuído a 12 cursos, o correspondente a 41,4 \% dos 29 cursos da categoria.
Dos demais cursos participantes, 11 ( 37,9 \% ) receberam conceito 2, e quatro ( 13,8 \% ) obtiveram conceito 4.
Os conceitos 1 e 5 foram obtidos por apenas um curso ( 3,4 \%, cada ).
Nenhum curso ficou Sem Conceito ( SC ).
Dos 65 cursos participantes na Modalidade Presencial, o conceito 4 foi o valor modal, atribuído a 28 cursos ( 43,1 \% ).
Dos demais cursos nessa Modalidade de Ensino, 21 cursos ( 32,3 \% ) receberam conceito 3, sete cursos ( 10,8 \% ) receberam conceitos 2 e 5, um curso ( 1,5 \% ) recebeu conceito 1, e um curso ( 1,5 \% ) ficou Sem Conceito ( SC ).
Dos 20 cursos participantes da Modalidade a Distância, o conceito modal foi 2, atribuído a 13 cursos ( 65,0 \% ).
Dos demais cursos nessa Modalidade, três cursos ( 15,0 \% ) receberem conceito 3, e dois ( 10,0 \% ) receberam conceito 1.
Os conceitos 4 e 5 foram obtidos por apenas um curso ( 5,0 \%, cada ).
Nessa Modalidade de Ensino, nenhum curso ficou Sem Conceito ( SC ).
Na região Norte, foi oferecido apenas um curso na Modalidade a Distância ( 100 \% ), o qual recebeu o conceito 1.
Dos cinco cursos que foram oferecidos na Modalidade Presencial, dois ( 40 \% ) receberam conceito 2, e os outros foram avaliados com conceitos 1, 3, e 4, cada um com um curso ( 20 \% ).
Todos os seis cursos oferecidos na região Norte são da rede pública com distribuição bimodal nos conceitos 1 e 2 com dois cursos cada ( 33,3 \% ).
Os dois cursos restantes foram alocados aos conceitos 3 e 4.
Na região Norte nenhum curso, independente da Categoria Administrativa ou da Modalidade de Ensino, ficou Sem Conceito ( SC ) ou recebeu o conceito 5 Na região Nordeste, a rede privada não teve cursos participantes, sendo assim, as IES Públicas concentraram 100 \% dos cursos correspondentes ao total dessa região.
Sete cursos ( 46,7 \% ) foram avaliados com conceito 4, o conceito modal.
Quatro cursos ( 26,7 \% ) obtiveram conceito 2, e outros quatro ( 26,7 \% ) obtiveram conceito 3.
Nessa combinação de categoria e região, nenhum curso ficou Sem Conceito ( SC ) ou com os demais conceitos.
A maioria dos cursos oferecidos na região Nordeste eram Presenciais ( 80 \% ), com a moda no conceito 4 com seis cursos.
No Ensino a Distância a moda ficou no conceito 2 com dois cursos.
Na região Sudeste, a proporção de cursos da rede privada foi de 57,6 \%, correspondendo a 19 dos 33 cursos participantes.
Nessa categoria e região, o conceito modal foi 3, atribuído a oito ( 42,1 \% ) dos 19 cursos.
Os demais foram avaliados com conceito 1 ( um curso ), conceito 2 ( sete cursos ), conceito 4 ( dois cursos ) e conceito 5 ( um curso ).
Dos 14 cursos oferecidos por Instituições Públicas na região Sudeste, os conceitos 4 e 5 foram atribuídos a seis cursos, cada, o valor considerado uma moda com duas classes contíguas.
Nessa região, dois cursos ( 14,3 \% ) receberam conceito 3, e nenhum curso recebeu os conceitos 1 ou 2 ou, ainda, ficou Sem Conceito ( SC ).
Também na região Sudeste, grande parte dos cursos foi oferecida na Modalidade Presencial : 24 de 33 cursos.
Para esses, a moda foi o conceito 3, com nove cursos.
Os demais cursos Presenciais da região foram alocados nos conceitos 2 ( um curso ), 4 ( oito cursos ) e 5 ( seis cursos ).
Nenhum curso recebeu conceito 1 ou ficou Sem Conceito ( SC ).
Na região Sudeste, seis cursos ( 66,7 \% ) oferecidos na Modalidade Educação a Distância obtiveram conceito 2, o conceito modal.
E os conceitos 1, 3 e 5 foram recebidos por apenas um curso, cada um correspondendo a 11,1 \%.
As Instituições Privadas concentraram nove dos 23 cursos participantes da região Sul, 39,1 \% do total regional.
Desses, quatro cursos obtiveram conceito 3, o conceito modal.
Nessa combinação de Categoria Administrativa e Grande Região, nenhum curso recebeu os conceitos 1 ou 5, três cursos receberam conceito 2, e dois receberam o conceito 4.
As Instituições Públicas na região Sul participaram com 14 cursos ( 60,9 \% dos cursos da região ), aos quais foram atribuídos os conceitos : 2 ( um curso ), 3 ( três cursos ), 4 ( oito cursos, modal ) e 5 ( um curso ).
Apenas um curso da região Sul oferecido por IES Pública ficou Sem Conceito ( SC ), e nenhum recebeu conceito 1.
Na região Centro-Oeste, sete dos oito cursos participantes eram de Instituições Públicas ( 87,5 \% em termos regionais ).
Desses, três cursos receberam conceito 4, o modal.
Dos demais cursos, dois receberam o conceito 2 e dois receberam o conceito 3.
Nenhum curso ficou Sem Conceito ( SC ) ou recebeu os demais conceitos nessa combinação de categoria e região.
Em relação a modalidade de ensino, 62,5 \% dos cursos eram Presenciais ( cinco cursos ) e 37,5 \% a distância ( três cursos ).
Para os cursos Presenciais, a moda foi estabelecida no conceito 4 com três cursos e os outros dois cursos receberam os conceitos 2 e 3.
Para os cursos a Distância, a moda foi estabelecida no conceito 2 com dois cursos e o outro curso recebeu o conceito 3.
No Gráfico 5.2, apresenta -se a distribuição do Conceito Enade, segundo a Categoria Administrativa da IES.
Os cursos em IES Públicas ( linha azul ) apresentam uma poligonal mais à direita do que os em IES Privadas ( linha verde ) e, consequentemente, uma distribuição de Conceito Enade com valores maiores. 5.3 CONCEITOS POR ORGANIZAÇÃO ACADÊMICA E GRANDE REGIÃO Na Tabela 5.3, encontra -se a distribuição dos conceitos atribuídos aos cursos participantes do Enade/2021, na Área de Artes Visuais ( Licenciatura ), por Organização Acadêmica, segundo as Grandes Regiões brasileiras.
Dos 85 cursos de Artes Visuais ( Licenciatura ) participantes, 68 eram oferecidos em Universidades, nove em Centros Universitários, quatro em Faculdades, e outros quatro, em CEFET/IF.
Essa distribuição corresponde a, respectivamente, 80,0 \%, 10,6 \%, 4,7 \% e 4,7 \% dos cursos.
De acordo com os dados apresentados, de todos os oito cursos avaliados com conceito 5, sete eram vinculados a Universidades.
Os cursos vinculados a esse tipo de Organização Acadêmica tiveram o conceito 4 como conceito modal, atribuído a 26 cursos ( 38,2 \% ).
Os demais cursos vinculados a Universidades receberam os conceitos 1 ( dois cursos ), 2 ( 13 cursos ) e 3 ( 19 cursos ).
Apenas um curso ficou Sem Conceito ( SC ).
Dos cursos em Centros Universitários, os conceitos 2 e 3 ( valores contíguos constituindo a moda ) foram atribuídos a três cursos ( 33,3 \%, cada ).
Nenhum curso vinculado a esse tipo de Organização Acadêmica ficou Sem Conceito ( SC ) ou recebeu conceito 5.
Os demais cursos receberam os conceitos 1 ( um curso ) e 4 ( dois cursos ).
Dos quatro cursos mantidos por Faculdades, dois ( 50,0 \% ) receberam conceito 3, o conceito modal.
Na sequência, os conceitos 2 e 5 foram atribuídos a apenas um curso cada ( 25,0 \% ).
Dos cursos oferecidos em Faculdades, nenhum curso vinculado a esse tipo de Organização Acadêmica ficou Sem Conceito ( SC ) ou recebeu os conceitos 1 ou 4.
Na Organização Acadêmica CEFET/IF, dos quatro cursos, três ( 75,0 \% ) receberam conceito 2, o modal.
E apenas um curso ( 25 \% ) recebeu o conceito 4.
Os demais conceitos não foram contemplados.
Considerando -se, separadamente, as regiões brasileiras, verifica -se que, na região Norte, as Universidades participaram com apenas seis dos 68 cursos avaliados neste tipo de Organização Acadêmica, aos quais foram atribuídos conceitos 1 e 2 ( dois cursos cada ), e conceitos 3 e 4 ( um curso cada ).
Nenhum curso ficou Sem Conceito ( SC ) ou recebeu conceito 5 nessa combinação de organização e região.
Os Centros Universitários, as Faculdades, e os CEFET/IF não foram representados na região Norte.
Na região Nordeste, as Universidades participaram com 13 dos 15 cursos da Área de Artes Visuais ( Licenciatura ) oferecidos.
Nenhum desses 13 cursos ficou Sem Conceito ( SC ).
O conceito modal 4 foi atribuído a seis cursos.
Quatro cursos obtiveram conceito 3 e outros três cursos, o conceito 2.
Nenhum curso oferecido por Universidades no Nordeste recebeu conceitos 1 ou 5.
Os Centros Universitários e as Faculdades não foram representados na região Nordeste.
E os CEFET/IF foram representados por apenas dois cursos, sendo que um desses cursos recebeu conceito 2, e o outro recebeu conceito 4.
Na região Sudeste, as Universidades concentraram 23 dos 33 cursos de Artes Visuais ( Licenciatura ) da região.
Dos cursos oferecidos em Universidades, nessa região, os conceitos 3, 4 e 5, atribuídos a seis cursos, cada um, constituindo uma moda dividida em três conceitos contíguos.
Cinco cursos receberam o conceito 2.
Nenhum curso ficou Sem Conceito ( SC ) ou recebeu o conceito 1.
Os Centros Universitários participaram com seis cursos na região Sudeste, dos quais quatro obtiveram os conceitos 3 e 4 ( moda em dois conceitos contíguos ), dois cursos para cada um desses conceitos.
Os conceitos 1 e 2 foram atribuídos a um curso, cada.
Nenhum curso ficou Sem Conceito ( SC ) nessa combinação de organização acadêmica e região ou recebeu o conceito 5.
Já as Faculdades foram representadas por quatro cursos nessa região, dos quais nenhum ficou Sem Conceito ( SC ) ou recebeu os conceitos 1 ou 4.
O conceito modal 3 foi atribuído a dois cursos, um curso recebeu conceito 5 e outro o conceito 2.
Os CEFET/IF não foram representados na região Sudeste.
Dos 23 cursos da região Sul, 19 eram mantidos por Universidades, para os quais o conceito modal foi 4, atribuído a dez cursos.
Os demais cursos receberam os seguintes conceitos : 3 ( seis cursos ), 5 ( um curso ) e 2 ( um curso ).
Nesse tipo de organização, apenas um curso da região Sul ficou Sem Conceito ( SC ), e nenhum curso recebeu o conceito 1.
Dos três cursos participantes de Centros Universitários da região Sul, nenhum ficou Sem Conceito ( SC ) ou recebeu os conceitos 1, 4 e 5.
Os conceitos restantes foram atribuídos a dois cursos, o conceito modal 2, e a um curso, o conceito 3.
As Faculdades não foram representadas na região Sul.
Os CEFET/IF participaram com somente um curso, ao qual foi atribuído o conceito 2.
Na região Centro-Oeste, sete dos oito cursos oferecidos eram mantidos por Universidades.
Para essa Organização Acadêmica, os conceitos 2 e 3 foram atribuídos a dois cursos, cada, e o conceito 4, o modal, foi atribuído a três cursos.
Nesse tipo de organização, nenhum curso recebeu os demais conceitos ou ficou Sem Conceito ( SC ).
Os Centros Universitários e Faculdades não foram representados na região Centro-Oeste, e os CEFET/IF participaram com somente um curso, ao qual foi atribuído conceito 2.
No Gráfico 5.3, apresenta -se a distribuição do Conceito Enade dos cursos de Artes Visuais ( Licenciatura ), segundo a Organização Acadêmica da IES.
Os cursos em Universidades ( linha azul ) apresentam uma poligonal mais à direita do que as demais, denotando uma melhor distribuição de Conceitos Enade.
No outro extremo, encontram -se os cursos mantidos por CEFET/IF ( linha \textbf{laranja} ), apresentando as respectivas poligonais mais à esquerda e, consequentemente, distribuição de Conceitos Enade com menores valores.
Os cursos mantidos por Centros Universitários ( linha verde ) apresentam poligonais mais centrais e, consequentemente, distribuição de Conceitos Enade com valores intermediários.

% matched lemmas: ajudar, laranja
\end{textsample}
